\printedpage{139}
% head: A. Grothendieck : Dérivateurs, Ch. XVIII, Édition des Universités de Montpellier II et Paris VII
% foot: Version 17 décembre 2025

\noindent\textbf{Corollary 4.13.39.} {\itshape Under the conditions of 4.13.38, suppose that $K$ contains an object $e$ maximal in $J$. Let $\Sigma$ be the set of endomorphisms $p$ of $e$ in $J$, such that for every double arrow $k \overset{u}{\underset{v}{\rightrightarrows}} e$ \textup{in} $K$, we have $pu=pv$. Then $\Sigma$ is stable under composition, and there exists a nonempty subset $\Sigma_\infty$ \edbr{[}$=\Sigma_0$\edbr{]} of $\Sigma$, such that for $p,p'\in\Sigma_0$, exists
\msp{page 367}
$q\in\Sigma_0$ such that $qp=qp'$ (i.e. such that the subcategory $J_0$ of $J$ whose set of objects is $\{e\}$, and set of arrows $\{\id_e\}\cup\Sigma_\infty$, is filtering). If $\Sigma_0$ is chosen thus, let $\mathcal{F}$ be the union of $\Fl K$ and of the set of composite arrows
\[
k \xrightarrow{\ u\ } e \xrightarrow{\ p\ } e \xrightarrow{\ v\ } k'\ , \tag*{$(*)$}
\]
with $u,v\in\Fl J$, $p\in\Sigma_\infty$. Then $\mathcal{F}$ is the set of arrows of a subcategory $K'$ of $J$, having the same objects as $J$, and admitting $e$ as greatest element. Moreover, $K'$ contains $K$ and is filtering.}

\noindent\textsc{Proof}. That $\Sigma$ is stable under composition is clear. To construct $\Sigma_\infty$ having the stated properties, we note that, $J$ being large relative to $\pi$, $\Sigma$ is nonempty (\footnote[162]{cf. argument below for construction of $p_i$.}). We construct by induction a sequence of elements
\[
p_0\ ,\ p_1\ ,\ \ldots\ ,\ p_n\ ,\ \ldots
\]
of $\Sigma$, taking for $p_0$ any element of $\Sigma$, and for $i\geqslant 1$, taking $p_i\in\Sigma$ such that for every double arrow $(p,p')$ of the submonoid $M_{i-1}$ of $\End_J(e)$ generated by $(p_j)_{0\leqslant j\leqslant i-1}$, we have $p_ip=p_ip'$. Such a $p_i$ exists, for taking the set of double arrows of the preceding form $(p,p')$, with $p,p'\in M_{i-1}$, and of those of the form $(u,v) : k \rightrightarrows e$, with $u,v\in\Fl J$, the set of
\msp{page 368}
all these double arrows has cardinality $\leqslant\pi$ ($M_{i-1}$ being countable, hence of cardinality $\leqslant\pi$), hence there exists an arrow $e\longrightarrow j$ in $J$ which equalizes them all ($J$ being large relative to $\pi$). Composing if necessary with $j\longrightarrow e$ ($e$ being a maximal object of $J$), we find the desired $p_i : e\longrightarrow e$.

Let then $\Sigma_0$ be the multiplicative subset of $\Sigma$ generated by the $p_i$ ($i\in\mathbf{N}$), direct limit of the similar sets $\Sigma_i$ generated by the $p_j$ for $j\in[0,i]$. If $p$, $p'$ are in $\Sigma_\infty$, they are in one and the same $\Sigma_i$, hence we will have $p_{i+1}p=p_{i+1}p'$, which proves that $\Sigma_\infty$ has the required filtering property. % [sic] il sont -> ils sont, agreement of the verb with the plural subject

Let us now prove that $\mathcal{F}$ is the set of arrows of a subcategory $K'$ of $J$, having $\Ob K$ as set of objects. It remains to establish stability under composition $vu$. But if $u$ and $v$ are in $\Fl K$, or if one of the two is in $\Fl K$, the other not (hence of the form $(*)$), this stability is evident. Remains to see the case where neither \edbr{[}$u$\edbr{]} nor $v$ is in $\Fl K$, hence both are of the form $(*)$:
\[
k \xrightarrow{\ u\ } e \xrightarrow{\ p\ } e \xrightarrow{\ v\ } k' \xrightarrow{\ u'\ } e \xrightarrow{\ p'\ } e \xrightarrow{\ v'\ } k''\ ,
\]
with $u,v,u',v'\in\Fl K$, $p,p'\in\Sigma_\infty$. Since $(u'v,\id_e)$ is a double arrow of
\printedpage{140}
% head: A. Grothendieck : Dérivateurs, Ch. XVIII, Édition des Universités de Montpellier II et Paris VII
% foot: Version 17 décembre 2025
\msp{page 369}
$K$ with target $e$, we have
\[
p'u'v \;=\; p'\id_e \;=\; p'\ ,
\]
hence the composite in question is also written
\[
k \xrightarrow{\ u\ } e \xrightarrow{\ p'p\ } e \xrightarrow{\ v'\ } k''\ ,
\]
where $p'p\in\Sigma_\infty$ since $\Sigma_\infty$ is stable under composition, hence it is in $\mathcal{F}$, q.e.d.

It is clear that $K'$ contains $K$, and having the same objects as $K$ which admits $e$ as greatest element, that $e$ is a greatest element of $K'$. Remains to prove that $K'$ is filtering. Since $K'$ has a greatest element, it remains to see that every double arrow $k\rightrightarrows l$ in $K'$ can be equalized in $K'$. Composing if necessary with $l\longrightarrow e$, we can suppose $l=e$. Now we see that every arrow $k\longrightarrow e$ in $K'$ is of the form
\[
k \xrightarrow{\ u\ } e \xrightarrow{\ p\ } e \xrightarrow{\ v\ } e\ ,
\]
with $u,v\in\Fl K$ and $p\in\Sigma_\infty\cup\{\id_e\}$. Composing if necessary with $p_0\in\Sigma_\infty$ fixed, this arrow becomes (since $p_0v=p_0\id_e=p_0$)
\[
k \xrightarrow{\ u\ } e \xrightarrow{\ p_0p\ } e\ , \qquad\quad \text{i.e. } p_0pu\ .
\]
For an arrow $v'p'u' : k\longrightarrow e$, we find $p_0p'u'$. Let then $q\in\Sigma_0$ be such that $q(p_0p)=q(p_0p')$, say $r$ ($q$ exists since $p_0p,p_0p'\in\Sigma_\infty$). Thus $(q(p_0pu),q(p_0p'u'))=(ru,ru')$, and we have $ru=ru'$ since $r\in\Sigma_\infty$. All told, $qp_0$ equalizes the two arrows $vpu$ and $v'p'u'$, q.e.d.

\msp{page 370}

\section*{4.14\quad Accessibility criteria for big categories}

\noindent\textbf{Theorem 4.14.1.} {\itshape Let $\mathcal{C}$ be a category equivalent to a small category, $\Phi$, $\Psi$ two small sets of categories, equivalent to small categories. For every category $I$, we denote by $I_{\bullet}$ (resp. ${}_{\bullet}I$) the right cone (resp. left) on $I$, obtained from $I$ by adding a final (resp. initial) object $e=e_I$. Let $\Lambda_1$ (resp. $\Lambda_2$) be a set of pairs $(I,\varphi)$ (resp. $(I,\psi)$), where $I\in\Phi$ (resp. $I\in\Psi$), and where
\[
\varphi\ :\ I \longrightarrow \mathcal{C} \qquad\quad \text{(resp. }\psi\ :\ I \longrightarrow \mathcal{C}\ \ \text{)} \tag*{\textup{(4.14.1.1)}}
\]
is a functor. For every category $\mathcal{N}$ stable (\footnote[163]{This stability property is not necessary for the definition of $\uHom^{\Lambda_2}_{\Lambda_1}(\mathcal{C},\mathcal{N})$, and it is almost certain that the theorem is valid without making this assumption on $\mathcal{N}$. Cf. 4.14.16, where we see that it is in any case unnecessary for $\Psi$.}), under $\varinjlim$ of type $\Phi$ and under $\varprojlim$ of type $\Psi$ we denote by
\[
\uHom^{\Lambda_2}_{\Lambda_1}(\mathcal{C},\mathcal{N})\ \subset\ \uHom(\mathcal{C},\mathcal{N}) \tag*{\textup{(4.14.1.2)}}
\]
the strictly full subcategory of $\uHom(\mathcal{C},\mathcal{N})$ formed of the functors $f:\mathcal{C}\longrightarrow\mathcal{N}$ having the following property:
\begin{enumerate}[leftmargin=2.5em]
\item[\textit{a)}] $\forall\ (I,\varphi)\in\Lambda_1$, the composite functor
\[
I_{\bullet} \xrightarrow{\ \varphi\ } \mathcal{C} \xrightarrow{\ f\ } \mathcal{N}
\]
is $I$-left exact, i.e. makes $(f\varphi)(e_I)$ a $\varinjlim\limits_{I}$ of $f\varphi|I$.
\printedpage{141}
% head: A. Grothendieck : Dérivateurs, Ch. XVIII, Édition des Universités de Montpellier II et Paris VII
% foot: Version 17 décembre 2025
\msp{page 371}
\item[\textit{b)}] For every $(I,\psi)\in\Lambda_2$, the composite functor
\[
{}_{\bullet}I \xrightarrow{\ \psi\ } \mathcal{C} \xrightarrow{\ f\ } \mathcal{N}
\]
is $I$-right exact, i.e. makes $(f\psi)(e_I)$ an inverse limit of $(f\psi)|I$.
\end{enumerate}
\noindent With these notations: if $\mathcal{N}$ is an accessible category, the same is true of $\uHom^{\Lambda_2}_{\Lambda_1}(\mathcal{C},\mathcal{N})$.}

\noindent\textbf{Comment.} We will use this theorem chiefly when we have $\Lambda_2=\emptyset$ or $\Lambda_1=\emptyset$, in which case we write $\uHom_{\Lambda_1}(\mathcal{C},\mathcal{N})$, resp. $\uHom^{\Lambda_2}(\mathcal{C},\mathcal{N})$, for the category $\uHom^{\Lambda_2}_{\Lambda_1}(\mathcal{C},\mathcal{N})$. We will apply the theorem chiefly when $\mathcal{C}$ itself is stable under $\varinjlim$ of type $\Phi$ and under $\varprojlim$ of type $\Psi$, and when $\Lambda_1$ (resp. $\Lambda_2$) is already formed of functors $\varphi$ (resp. $\psi$) which are $I$-left exact (resp. right exact). Thus $\uHom^{\Lambda_2}_{\Lambda_1}(\mathcal{C},\mathcal{N})$ is then formed of the $f:\mathcal{C}\longrightarrow\mathcal{N}$ which commute with certain $\varinjlim$ of type $\Phi$, and with certain $\varprojlim$ of type $\Psi$. In particular, we find:

\msp{page 372}

\noindent\textbf{Corollary 4.14.2.} {\itshape Let $\mathcal{C}$ be a category equivalent to a small category, $\mathcal{N}$ a $\mathfrak{U}$-category, $\Phi$ and $\Psi$ two small sets of categories $I$, equivalent to small categories. We suppose $\mathcal{C}$ and $\mathcal{N}$ stable under $\varinjlim$ of type $\Phi$, under $\varprojlim$ of type $\Psi$. Let $\uHom^{\Psi}_{\Phi}(\mathcal{C},\mathcal{N})$ be the strictly full subcategory of $\uHom(\mathcal{C},\mathcal{N})$ formed of the $f:\mathcal{C}\longrightarrow\mathcal{N}$ which commute with $\varinjlim$ of type $\Phi$ and with $\varprojlim$ of type $\Psi$. Then, if $\mathcal{N}$ is accessible, the same is true of $\uHom^{\Psi}_{\Phi}(\mathcal{C},\mathcal{N})$.}

Here are some remarkable special cases of this corollary:

\noindent\textbf{Corollary 4.14.3.} {\itshape Let $\mathcal{M}$, $\mathcal{N}$ be two accessible categories, and $\pi_0$ an infinite cardinal, such that $\mathcal{M}$ satisfies $L_{\pi_0}$. Let $\uHom_{\pi_0}(\mathcal{M},\mathcal{N})$ be the strictly full subcategory of $\uHom(\mathcal{M},\mathcal{N})$ formed of the $\pi_0$-accessible functors. Then $\uHom_{\pi_0}(\mathcal{M},\mathcal{N})$ is accessible.}

\msp{page 373}

(The case where $\pi_0$ is useful for $\mathcal{M}$ has already been treated.)

\noindent\textsc{Proof}. In the general case, let $\pi>\pi_0$ be a useful cardinal for $\mathcal{M}$, so that we have
\[
\mathcal{M} \xleftarrow{\ \sim\ } \Ind_\pi(\mathcal{C}) \qquad\quad \text{with } \mathcal{C}=\mathcal{M}_\pi\ .
\]
We can suppose $\mathcal{M}=\Ind_\pi(\mathcal{C})$. The fact that $\mathcal{M}$ satisfies $\mathrm{L}_{\pi_0}$ implies that $\mathcal{C}$ is stable under $\varinjlim$ large relative to $\pi_0$, and of cardinality $\leqslant\pi$. Let $\Phi$ be the set of categories of this nature. $\Phi$ is not small, but there exists a small $\Phi_0\subset\Phi$, such that every $I\in\Phi$ is isomorphic to an $I_0\in\Phi_0$, which will suffice for the validity of 4.14.1 and 4.14.2 (where here we will put $\Psi=\emptyset$). We can consider $\uHom_{!\,\pi_0}(\mathcal{M},\mathcal{N})$ as a strictly full subcategory of $\uHom_{!\,\pi}(\mathcal{M},\mathcal{N})\xrightarrow{\ \sim\ }\uHom(\mathcal{C},\mathcal{N})$ (cf. 4.12.3, p. 275). For a functor $f:\mathcal{M}\longrightarrow\mathcal{N}$ which is $\pi$-accessible, to say that it is even $\pi_0$-accessible amounts to saying that its restriction to $\mathcal{C}$, i.e. its image in $\uHom(\mathcal{C},\mathcal{N})$, commutes with $\varinjlim$ of type $\Phi$. (Cf. 4.13.20, page 341 (\footnote[164]{The reference is not quite adequate. We should gather in a separate paragraph all the statements of this type, some of which use the indicial representation theorem.}), where we take $\alpha=\alpha_\pi$, noting that $J_{\alpha_\pi}$ large relative to $\pi_0$ since $\pi>\pi_0$).

\msp{page 374}

Thus we find
\[
\uHom_{!\,\pi_0}(\mathcal{M},\mathcal{N}) \xrightarrow{\ \sim\ } \uHom_{\Phi}(\mathcal{C},\mathcal{N})\ ,
\]
\printedpage{142}
% head: A. Grothendieck : Dérivateurs, Ch. XVIII, Édition des Universités de Montpellier II et Paris VII
% foot: Version 17 décembre 2025
and corollary 4.14.2 gives the desired conclusion.

\noindent\textbf{Corollary 4.14.4.} {\itshape Let $\mathcal{M}$, $\mathcal{N}$ be two accessible categories, $\Phi$ and $\Psi$ two sets of categories as in 4.14.2, we suppose that $\mathcal{M}$ and $\mathcal{N}$ satisfy $L_{\pi_0}$, and \edbr{[are]} stable under $\varinjlim$ of type $\Phi$ and under $\varprojlim$ of type $\Psi$. Let $\uHom^{\Psi}_{\Phi,\pi_0}(\mathcal{M},\mathcal{N})$ be the strictly full subcategory of $\uHom(\mathcal{M},\mathcal{N})$ formed of the $\pi_0$-accessible functors which commute with $\varinjlim$ of type $\Phi$, and with $\varprojlim$ of type $\Psi$. Then $\uHom^{\Psi}_{\Phi,\pi_0}(\mathcal{M},\mathcal{N})$ is accessible.}

\noindent\textsc{Proof}. This time, we choose $\pi>\pi_0$ such that we have $\cardred I\leqslant\pi$ for every $I$ in $\Phi$ or in $\Psi$. The category under consideration is then interpreted (putting again $\mathcal{C}=\mathcal{M}_\pi$) as
\msp{page 375}
\[
\uHom^{\Psi}_{\Phi,\pi_0}(\mathcal{M},\mathcal{N}) \xrightarrow{\ \sim\ } \uHom^{\Psi}_{\Phi'}(\mathcal{C},\mathcal{N})\ ,
\]
where $\Phi'$ is formed of the categories which are either in $\Phi$, or which are large relative to $\pi_0$ and of cardinality $\leqslant\pi$. It therefore suffices to apply 4.14.2 once more.

\noindent\textbf{Corollary 4.14.5.} {\itshape Let us consider the following three sets:
\begin{enumerate}[leftmargin=2.5em]
\item[\textit{a)}] $\Phi_1$ is formed of all categories equivalent to small categories.
\item[\textit{b)}] $\Phi_2$ as a), but restricting to filtering categories.
\item[\textit{c)}] $\Phi_3$ formed of the discrete categories, and of those which are large relative to $\pi_0$ ($\pi_0$ a given infinite cardinal).
\end{enumerate}
\noindent Let $\Phi$ be one of the sets of categories $\Phi_1$, $\Phi_2$, $\Phi_3$, $\Phi_2\cup\Phi_3$ or the union of one of these and a small set $\Phi_0$ of essentially small categories (\footnote[165]{We could give here a statement which includes 4.14.4 by adding a
\begin{enumerate}[leftmargin=2.5em]
\item[d)] $\Phi_4$ formed of the categories large relative to $\pi_0$,
\end{enumerate}
\noindent and take $\Phi$ of the form of $\Phi_1$, $\Phi_2$, $\Phi_4$, $\Phi_3$, $\Phi_3\cup\Phi_2$, $\Phi_3\cup\Phi_4$, or the union of these and a $\Phi_0$.}). Let $\Psi$ be a small set of categories. Let $\mathcal{M}$, $\mathcal{N}$ be accessible categories, stable under $\varinjlim$ of type $\Phi$, and under $\varprojlim$ of type $\Psi$. Then $\uHom^{\Psi}_{\Phi}(\mathcal{M},\mathcal{N})$ is accessible.}

\msp{page 376}

\noindent\textsc{Proof}. Let $\pi'_0\geqslant\pi_0$ be a useful cardinal for $\mathcal{M}$, $\pi$ a cardinal having the following properties.
\begin{enumerate}[leftmargin=2.5em]
\item[a)] $\pi\geqslant\sup(\pi'_0,\cardred\mathcal{M}_{\pi'_0})$, $\pi>\pi_0$.
\item[b)] $\pi^{\pi'_0}=\pi$.
\item[c)] $I\in\Phi_0\cup\Psi\Longrightarrow\cardred I\leqslant\pi$.
\end{enumerate}
\noindent Let $\mathcal{C}=\mathcal{M}_\pi$, whence $\mathcal{M}\simeq\Ind_\pi(\mathcal{C})$. In all the cases under consideration, $\Phi$ contains the set of essentially small categories which are large relative to $\pi_0$, a fortiori those which are large relative to $\pi$, hence the category $\uHom^{\Psi}_{\Phi}$ is interpreted as a strictly full subcategory of
\[
\uHom(\mathcal{C},\mathcal{N})\ .
\]
If on the other hand $\Phi_\pi$ denotes the set of $I\in\Phi$ such that $\card I\leqslant\pi$, we find
\[
\uHom^{\Psi}_{\Phi}(\mathcal{M},\mathcal{N})\ \simeq\ \uHom^{\Psi}_{\Phi_\pi}(\mathcal{C},\mathcal{N})\ ,
\]
and we apply 4.14.2.

\printedpage{143}
% head: A. Grothendieck : Dérivateurs, Ch. XVIII, Édition des Universités de Montpellier II et Paris VII
% foot: Version 17 décembre 2025
\msp{page 377}

In particular, we find:

\noindent\textbf{Corollary 4.14.6.} {\itshape Let $\mathcal{M}$, $\mathcal{N}$ be two accessible categories, suppose $\mathcal{M}$ stable under small $\varinjlim$ (i.e. $\mathcal{M}$ is a paratopos category). Then the category $\uHom_{!}(\mathcal{M},\mathcal{N})$ of functors $f:\mathcal{M}\longrightarrow\mathcal{N}$ which commute with small $\varinjlim$ is accessible. If $\mathcal{N}$ is itself also a paratopos, $\uHom_{!}(\mathcal{M},\mathcal{N})$ is a paratopos.}

\noindent\textsc{Proof}. The first assertion is the special case of 4.14.5 corresponding to the case where $\Phi=\Phi_0$, $\Psi=\emptyset$. The second assertion follows from it, since $\uHom_{!}(\mathcal{M},\mathcal{N})$ is visibly stable under small $\varinjlim$ , hence it is a paratopos being accessible (4.8.9, page 207).

\noindent\textbf{Corollary 4.14.7.} {\itshape Let $\Phi$, $\Psi$, $\mathcal{M}$ be as in 4.14.5. Then for every accessible category $\mathcal{N}$, the category
\msp{page 378}
\[
\uHom^{\Phi}_{\Psi}(\mathcal{M}^{\circ},\mathcal{N})\ \simeq\ (\uHom^{\Psi}_{\Phi}(\mathcal{M},\mathcal{N}^{\circ}))^{\circ}
\]
is accessible.}

\noindent\textsc{Proof}. We proceed as for 4.14.5, to make explicit
\[
\uHom^{\Psi}_{\Phi}(\mathcal{M},\mathcal{N}^{\circ}) \xrightarrow{\ \sim\ } \uHom^{\Psi}_{\Phi_\pi}(\mathcal{C},\mathcal{N}^{\circ})\ ,
\]
whence
\[
\uHom^{\Phi}_{\Psi}(\mathcal{M}^{\circ},\mathcal{N}) \xrightarrow{\ \sim\ } \uHom^{\Phi_\pi}_{\Psi}(\mathcal{C}^{\circ},\mathcal{N})\ ,
\]
and we apply 4.14.2 to $\mathcal{C}^{\circ}$, $\Psi$, $\Phi_\pi$.

\noindent\textbf{Corollary 4.14.8.} {\itshape Let $\mathcal{M}$ and $\mathcal{N}$ be two paratopos categories. Then
\[
\mathcal{M}\boxtimes\mathcal{N}\ \simeq\ \uHom^{!}(\mathcal{M}^{\circ},\mathcal{N}) \qquad\quad \text{(cf. 4.2 and 4.3)}
\]
is a paratopos.}

Indeed, by virtue of 4.14.7 (with $\Phi=\Phi_\pi$ [?], $\Psi=\emptyset$), it is an accessible category.

\msp{page 379}

On the other hand, it is immediate that it is stable under small $\varprojlim$ , $\mathcal{N}$ being so. Thus we conclude again by 4.8.9 (page 207) that it is a paratopos category.

\noindent\textbf{Corollary 4.14.9.} {\itshape Let $\mathcal{M}$, $\mathcal{N}$ be two paratopos categories, and $\uHom^{!}(\mathcal{M},\mathcal{N})$ the category of functors $\mathcal{M}\longrightarrow\mathcal{N}$ which commute with small $\varprojlim$ . We then have an equivalence of categories
\[
\begin{aligned}
\uHom^{!}(\mathcal{M},\mathcal{N}) &\ \xrightarrow{\ \sim\ }\ \uHom_{!}(\mathcal{N},\mathcal{M})^{\circ}\\
f &\ \longmapsto\ \text{left adjoint of } f\ ,
\end{aligned}
\]
\noindent\textup{\textbf{!}}\ \textup{\textbf{!}} and consequently the category \textup{\textbf{opposite}} to $\uHom^{!}(\mathcal{M},\mathcal{N})$ is accessible.} % Verified against PDF leaf 25: a nonverbal marginal curved mark; see the correction register.

This follows from the

\noindent\textbf{Lemma 4.14.10} (for the record). {\itshape Let $\mathcal{M}$ be \edbr{[a]} pseudo-topos category, $f:\mathcal{M}\longrightarrow\mathcal{N}$ a functor. For $f$ to admit a left adjoint (resp. a right adjoint), it is necessary and sufficient that $f$ commute with small $\varprojlim$ (resp. with small $\varinjlim$ ).}